%--------------------------------------------------------------------------
%  ¿Importa conocer la BEMF con exactitud?  — reunión con el asesor, 2026-10-23
%  Figuras: python3 python_scripts/make_advisor_figs.py  → fig/asesor_*.pdf
%  Guión:   presentation/resultados_asesor_guion.md
%  Compilar: cd presentation && tectonic -X compile resultados_asesor.tex
%--------------------------------------------------------------------------
\documentclass[aspectratio=169,xcolor=dvipsnames]{beamer}
\makeatletter
\def\input@path{{theme/}}
\makeatother
\usetheme{SimplePlus}
\usepackage{graphicx}
\usepackage{booktabs}
\usepackage{amsmath,amssymb}
\graphicspath{{fig/}}
\setbeamerfont{frametitle}{size=\large,series=\bfseries,parent=structure}

\newcommand{\figfull}[1]{\begin{center}\includegraphics[width=\linewidth]{#1}\end{center}\vspace{-0.5cm}}
\newenvironment{notas}{\footnotesize\begin{itemize}\setlength{\itemsep}{1pt}\setlength{\parskip}{0pt}}{\end{itemize}}

\title{¿Importa conocer la BEMF con exactitud?}
\subtitle{Lo que encontramos al cerrar el lazo con el ADALINE en el motor real}
\author{Alejandro León}
\institute{Instituto de Ingeniería \\ Universidad Nacional Autónoma de México}
\date{23 de octubre de 2026}

\begin{document}

\begin{frame}
  \titlepage
\end{frame}

%--------------------------------------------------------------------------
\begin{frame}{Qué queríamos comprobar}
  \begin{block}{La idea de la tesis}
    Si el controlador predictivo aprende la forma real de la BEMF mientras el motor gira,
    su predicción es mejor y el rizo de corriente y de par baja.
  \end{block}
  \vspace{0.15cm}
  En simulación esto funcionó. La pregunta de este semestre era si también funciona en el motor real.
  Para responderla hicimos cuatro cosas:
  \begin{enumerate}\setlength{\itemsep}{1pt}
    \item Medimos la BEMF del motor con el inversor apagado, para tener una referencia que no dependa del controlador.
    \item Corrimos el control con tres modelos de BEMF: el nominal, uno equivocado a propósito y el que aprende el ADALINE.
    \item Comparamos no sólo el error RMS, sino también el valor medio de la corriente y sus armónicos.
    \item Probamos qué tan estable es el aprendizaje en línea con distintas ganancias.
  \end{enumerate}
  \vspace{0.15cm}
  La respuesta corta: en este motor, el modelo de BEMF casi no influye en el resultado del control,
  y eso tiene una explicación numérica clara.
\end{frame}

%--------------------------------------------------------------------------
\begin{frame}{La BEMF del motor es prácticamente una sinusoide}
  \figfull{asesor_bemf_banco.pdf}
  \begin{notas}
    \item La curva medida y la sinusoide ideal caen una sobre la otra; la diferencia no pasa del $0.6\,\%$ del pico.
          La distorsión armónica total es de $0.4$ a $0.7\,\%$ y sólo el tercer y el quinto armónico se distinguen del ruido.
    \item Para comparar: el motor de la simulación tiene $2.6\,\%$ de distorsión y el modelo trapezoidal con el que se comparaba, $4.6\,\%$.
          Aquí hay muy poco que aprender.
  \end{notas}
\end{frame}

%--------------------------------------------------------------------------
\begin{frame}{Cuánto pesa un error de BEMF dentro del predictor}
  \footnotesize
  El predictor avanza la corriente un paso con la tensión aplicada menos la BEMF,
  $i(k{+}1)=i(k)+G\,[u(k)-e(k)-R\,i(k)]$ con $G=T_s/L=0.023$ A/V. Como se compensa un retardo, la predicción
  abarca dos pasos: un error $\delta e$ en la BEMF mueve la corriente predicha $2G\,\delta e$. Eso permite poner todo en una misma escala:
  \vspace{-0.1cm}
  \figfull{asesor_escalas.pdf}
  \begin{notas}
    \item La distorsión real de la BEMF vale $0.2$ mA en la predicción y un modelo equivocado un $25\,\%$, $9$ mA.
          El controlador trabaja con pasos de $19$ mA y el error que medimos es de $48$ mA.
  \end{notas}
\end{frame}

%--------------------------------------------------------------------------
\begin{frame}{Lo que predice esa cuenta y lo que medimos}
  \figfull{asesor_sensibilidad.pdf}
  \begin{notas}
    \item Si el error de modelo y el piso del controlador son independientes, el RMS es la suma cuadrática de los dos.
          La curva azul es esa predicción, sin ningún parámetro ajustado: es plana hasta errores de varias decenas de por ciento.
          Para subir el RMS un $10\,\%$ habría que equivocarse un $60\,\%$; quitando la BEMF por completo llegaría a $60$ mA.
    \item Los dos puntos medidos caen sobre la parte plana: $48\pm2.5$ mA con el modelo nominal y $45\pm2$ mA con el equivocado un $25\,\%$.
  \end{notas}
\end{frame}

%--------------------------------------------------------------------------
\begin{frame}{En el experimento, el error RMS no distingue un modelo de otro}
  \figfull{asesor_rms_experimento.pdf}
  \begin{notas}
    \item Cuatro corridas a la misma velocidad y en la misma sesión: el modelo nominal da $48$ mA, el equivocado $45$ y el aprendido $45$.
          Esas diferencias están dentro de lo que varía una corrida de otra.
    \item El error en el eje $d$ es casi igual al del eje $q$: eso apunta a un piso de ruido y de resolución, no a algo con la estructura de la BEMF.
    \item Cuando el ADALINE parte de un modelo equivocado, el error sube con el tiempo hasta $52\pm4$ mA (lámina de estabilidad).
  \end{notas}
\end{frame}

%--------------------------------------------------------------------------
\begin{frame}{Donde sí se nota el modelo: en el valor medio de la corriente}
  \figfull{asesor_sesgo_dc.pdf}
  \begin{notas}
    \item El valor medio se mide con $0.1$ mA de precisión, 25 veces mejor que el RMS, y ahí el modelo equivocado sí deja huella:
          desplaza la corriente media $12.6$ mA en el eje $d$ y $10.3$ mA en el eje $q$.
    \item La misma cuenta de dos pasos, $-2G\,\delta e$, predice $-5.0$ y $-7.6$ mA: el signo y el orden de magnitud correctos sin ajustar nada.
          Un error en el fundamental de la BEMF se vuelve un sesgo constante, no rizo.
  \end{notas}
\end{frame}

%--------------------------------------------------------------------------
\begin{frame}{El rizo que queda no viene de la forma de la BEMF}
  \figfull{asesor_armonicos.pdf}
  \begin{notas}
    \item Descomposición de la corriente por armónicos del ángulo eléctrico (26\,000 ventanas, resolución $0.12$ mA).
          El rizo de par aparece en el sexto armónico: $2$ mA en $i_q$ y $15$ mA en $i_d$.
    \item Con la distorsión medida en la BEMF, su forma explica como mucho $0.2$ mA de ese armónico: una décima del rizo en $i_q$
          y una centésima en $i_d$. El resto lo ponen el lazo y el inversor.
  \end{notas}
\end{frame}

%--------------------------------------------------------------------------
\begin{frame}{El ADALINE aprende la amplitud, pero arrastra un sesgo de fase}
  \figfull{asesor_identificacion.pdf}
  \begin{notas}
    \item Arrancando con una BEMF un $20\,\%$ más chica de lo real, el ADALINE la corrige en unos $110$ s y termina a $3\,\%$
          del valor medido con el inversor apagado.
    \item La fase no la corrige: arranque donde arranque termina entre $10$ y $12^\circ$ desplazada, y aparece un desbalance
          entre $\alpha$ y $\beta$ del $6$ al $8\,\%$. Sospechamos del observador, que usa la tensión pedida al inversor y no la aplicada;
          no está comprobado.
  \end{notas}
\end{frame}

%--------------------------------------------------------------------------
\begin{frame}{La adaptación en línea sólo es segura con ganancia pequeña y buena señal}
  \figfull{asesor_estabilidad.pdf}
  \begin{notas}
    \item Con la ganancia de la simulación el lazo se vuelve inestable: el error se triplica y salta la protección de sobrecorriente.
    \item Con una ganancia 20 veces menor y velocidad constante no estorba; con velocidad variable y BEMF pequeña, el ruido
          del observador domina y el error sube más del doble. Partiendo de un modelo equivocado sube despacio, un $11\,\%$ al final;
          la sensibilidad daría apenas $1$ mA y esto aún no lo entendemos.
  \end{notas}
\end{frame}

%--------------------------------------------------------------------------
\begin{frame}{Qué concluyo y qué quisiera decidir hoy}
  \footnotesize
  \textbf{Lo que queda establecido.} La BEMF de este motor es sinusoidal con menos de $0.7\,\%$ de distorsión.
  El error de corriente tiene un piso de $48$ mA que ponen el controlador y la medición, y un error de modelo de hasta
  un $25\,\%$ no lo mueve. El ADALINE identifica la amplitud en lazo cerrado.

  \vspace{0.1cm}
  \textbf{Lo que queda acotado.} Aprender la forma de la BEMF sólo paga cuando el error de modelo es comparable al piso
  del controlador. En este motor estamos más de un orden de magnitud por debajo de ese punto. La hipótesis no se refuta;
  se delimita.

  \vspace{0.1cm}
  \textbf{Lo que no sabemos.} De dónde sale el sesgo de fase del observador, y por qué el aprendizaje degrada el control
  más de lo que la sensibilidad predice. Todo se midió a 300 rpm y casi sin carga.

  \vspace{0.15cm}
  \begin{alertblock}{Opciones para cerrar la tesis}
    \begin{enumerate}\setlength{\itemsep}{1pt}
      \item Cerrar con este resultado: BEMF medida, identificador validado y la frontera cuantitativa de cuándo importa el modelo.
      \item Cambiar la métrica a lo que sí responde al modelo (valor medio y armónicos) y repetir con carga.
      \item Forzar un desajuste grande, emulando una BEMF distorsionada, para comprobar la curva de sensibilidad.
      \item Corregir el observador antes de seguir adaptando en línea.
    \end{enumerate}
  \end{alertblock}
\end{frame}

\end{document}
